\chapter{Fundamentação Teórica} \label{Cap:Fund_Teo}

\label{chp:fundamentacao}

\section{Considerações iniciais}
\label{sec:fund_consideracoes_iniciais}

A concepção de uma ferramenta robusta para a interpretação e categorização de faturas financeiras requer o domínio integrado de fundamentos oriundos de múltiplas subáreas da Inteligência Artificial e da Engenharia de Software. Este capítulo estabelece as bases conceituais estritamente instrumentais ao projeto, organizando-se em quatro pilares fundamentais: (i) mineração de textos e processamento de linguagem natural (Seção~\ref{sec:fund_pln}); (ii) visão computacional, análise de leiaute e reconhecimento óptico de caracteres (Seção~\ref{sec:fund_visao_ocr}); (iii) arquiteturas neurais profundas e agentes inteligentes baseados em grafos (Seção~\ref{sec:fund_arquiteturas_agentes}); e (iv) engenharia de dados, MLOps e sustentabilidade de modelos em produção (Seção~\ref{sec:fund_mlops}). A Seção~\ref{sec:fund_consideracoes_finais} encerra o capítulo.

\section{Mineração de textos e processamento de linguagem natural}
\label{sec:fund_pln}

O Processamento de Linguagem Natural (PLN) constitui o ramo da Inteligência Artificial dedicado a capacitar sistemas computacionais a compreender, interpretar, manipular e gerar linguagens humanas \cite{BPLN_livro:2024, chowdhury2003natural}. No domínio financeiro, o desafio central da mineração textual reside na natureza não estruturada, ruidosa e despadronizada das descrições de transações bancárias e de cartão de crédito, que exige tratamento com conhecimento específico do domínio em vez de abordagens textuais genéricas \cite{ta2023specialized}.

\subsection{Pré-processamento léxico e morfológico}
A primeira barreira técnica na modelagem de textos curtos é a dispersão léxica, atacada pelas operações determinísticas usuais de padronização (caixa baixa, normalização de acentos e supressão de pontuação por expressões regulares) e pela remoção de \textit{stopwords} \cite{BPLN_livro:2024, matos2010pre}. Contudo, listas genéricas universais mostram-se insuficientes no domínio financeiro. É indispensável incorporar heurísticas específicas do negócio para remover sufixos de naturezas jurídicas e corporativas (como ``LTDA'', ``S/A'', ``ME'', ``EPP'', ``EIRELI'') e prefixos de liquidação financeira (``COMPRA'', ``PARCELA'', ``DEBITO''), os quais atuam como ruídos que distorcem as distribuições estatísticas \cite{ta2023specialized}.

\subsection{Estratégias de tokenização: da palavra ao nível de byte}
\label{sec:fund_tokenizacao}
A tokenização segmenta a cadeia textual contínua em unidades atômicas de processamento denominadas \textit{tokens} \cite{BPLN_livro:2024}. Modelos tradicionais baseados em palavras inteiras (\textit{word-level}) dependem de vocabulários fechados predefinidos. Diante de faturas financeiras, em que erros de digitação, truncamentos e aglutinações são onipresentes, essa abordagem falha sistematicamente devido à explosão de termos fora do vocabulário (\textit{Out-Of-Vocabulary} -- OOV).

Para superar essa fragilidade, o estado da arte avançou para a tokenização em subpalavras (\textit{subwords}, como Byte-Pair Encoding e WordPiece) e, no limite da granularidade, para a tokenização no nível de caractere ou byte \cite{zhang2015character, xue2021byt5}. Modelos que operam diretamente sobre sequências de bytes UTF-8, a exemplo do ByT5 \cite{xue2021byt5}, eliminam completamente o vocabulário fixo. Essa característica confere ao modelo maior tolerância a erros ortográficos e capacidade de generalizar sobre siglas desconhecidas e nomes comerciais aglutinados (ex.: ``MERCADOPAG*LOJA123'').

O custo dessa granularidade é o comprimento das sequências: uma descrição de 30 caracteres, que um tokenizador de subpalavras representa com 8 a 12 unidades, torna-se uma sequência de 30 a 40 bytes (caracteres acentuados ocupam dois bytes em UTF-8), o que multiplica o custo da atenção -- quadrático no comprimento -- e exige mais exemplos para que o modelo aprenda composições de bytes que em subpalavras já viriam pré-compostas. É por isso que modelos em nível de byte tendem a exigir \textit{corpora} maiores para atingir o mesmo desempenho, uma consideração central na leitura dos resultados do ajuste fino do ByT5 no Capítulo~\ref{Cap:Ava_Exp}. Os classificadores esparsos desta pesquisa exploram a mesma intuição de granularidade por outra via: $n$-gramas de caractere (2 a 4) como atributos do TF-IDF, que capturam a morfologia recorrente das descrições sem depender de vocabulário fechado nem de treinamento neural.

\subsection{Vetorização estatística versus representações densas (\textit{embeddings})}
Para que algoritmos computacionais processem dados textuais, é imperativo convertê-los em representações numéricas vetoriais.

\subsubsection{TF-IDF (Term Frequency-Inverse Document Frequency)}
A técnica clássica de TF-IDF quantifica a importância de um termo $t$ em um documento $d$ pertencente a um corpus $D$, calculada pelo produto da frequência do termo ($TF$) pelo inverso da frequência no documento ($IDF$) -- noção de especificidade de termos introduzida por \citeonline{sparckjones1972statistical} e consolidada como esquema de ponderação por \citeonline{salton1988term} --, aqui na formulação de \citeonline{rajaraman2011mining}:
\begin{equation}
	\text{TF-IDF}(t, d, D) = \text{TF}(t, d) \times \log\left(\frac{|D|}{1 + |\{d \in D : t \in d\}|}\right)
\end{equation}
Embora gere matrizes esparsas de alta dimensionalidade que desconsideram o contexto e a ordem sintagmática das palavras, o TF-IDF permanece relevante como uma linha de base científica (\textit{baseline}) de referência rápida, interpretável e de baixo custo computacional \cite{rajaraman2011mining}.

\subsubsection{Embeddings densos e bancos de dados vetoriais}
\label{sec:fund_embeddings_rag}
A superação das limitações esparsas consolidou-se com a semântica distribucional profunda, cuja formulação moderna remonta aos vetores de palavras aprendidos por redes neurais rasas a partir do contexto de ocorrência (\textit{word2vec}) \cite{mikolov2013efficient}. Modelos como o Sentence-BERT (SBERT) \cite{reimers2019sentence}, implementados por meio de arquiteturas siamesas de Sentence Transformers (ex.: \modelemb{}), mapeiam sentenças inteiras em vetores densos e contínuos de dimensionalidade fixa (ex.: 384 dimensões).

Nesse espaço vetorial contínuo, a proximidade semântica entre duas transações $u$ e $v$ é mensurada diretamente pela similaridade de cosseno:
\begin{equation}
	\text{Simil}(u, v) = \frac{u \cdot v}{\|u\|_2 \|v\|_2} = \frac{\sum_{i=1}^{n} u_i v_i}{\sqrt{\sum_{i=1}^{n} u_i^2} \sqrt{\sum_{i=1}^{n} v_i^2}}
\end{equation}
Para viabilizar consultas em milissegundos sobre milhões de registros históricos, utilizam-se extensões vetoriais integradas ao banco de dados relacional, destacando-se a extensão pgvector para PostgreSQL. Essa infraestrutura é o substrato de dois usos distintos, que convém não confundir. O primeiro é a \textbf{busca vetorial por similaridade} propriamente dita: os vizinhos mais próximos de uma consulta são recuperados e a categoria é decidida a partir deles (por exemplo, por votação ponderada pela similaridade). O segundo é a \textbf{Geração Aumentada por Recuperação} (\textit{Retrieval-Augmented Generation} -- RAG) \cite{lewis2020retrieval}, em que os exemplos recuperados são \textit{injetados no prompt} de um modelo de linguagem, que gera a resposta condicionado a essa evidência -- o mecanismo pelo qual um modelo passa a decidir ``em cima'' de uma base de conhecimento externa, em regime \textit{few-shot} dinâmico. A solução desta pesquisa implementa o primeiro uso, como ramo autônomo de classificação e como critério de roteamento do agente; os vizinhos recuperados \textbf{não são repassados} ao modelo de linguagem, que opera em regime \textit{zero-shot} (Seção~\ref{sec:sol_agente_inteligente}). Por essa razão, a arquitetura correspondente é denominada ao longo do texto ``busca vetorial densa'', reservando-se o termo RAG ao conceito aqui definido e à discussão, no Capítulo~\ref{Cap:Conclusoes}, da injeção de vizinhos como evolução natural do ramo de raciocínio.

A recuperação eficiente nesse espaço vetorial de alta dimensionalidade não precisa ser exaustiva: extensões como o pgvector empregam estruturas de indexação por Busca Aproximada de Vizinhos Mais Próximos (\textit{Approximate Nearest Neighbor} -- ANN), que trocam a garantia de exatidão por ganhos expressivos de velocidade \cite{jegou2011product}. Na variante \textit{IVFFlat} (\textit{Inverted File with Flat compression}), o espaço vetorial é particionado em um número fixo de células (\textit{listas}) por meio de agrupamento por $k$-médias, e a consulta é comparada apenas contra os vetores das células cujos centroides lhe são mais próximos; na variante \textit{HNSW} (\textit{Hierarchical Navigable Small World}), constrói-se um grafo de vizinhança em múltiplas camadas, navegado de forma gulosa da camada mais esparsa à mais densa, com parâmetros que regem o grau do grafo ($m$) e o esforço de construção (\textit{ef\_construction}). Em ambas, há um compromisso entre velocidade e revocação: quanto mais agressiva a aproximação, maior a probabilidade de omitir o vizinho verdadeiramente mais próximo. Nesta pesquisa, o esquema físico cria um índice HNSW ($m = 16$, \textit{ef\_construction} $= 64$) sobre a tabela relacional de estabelecimentos (Apêndice~\ref{apendice:esquema_banco}); a coleção vetorial consultada pelo agente, mantida pela integração LangChain, não recebe índice aproximado e é percorrida por busca exaustiva -- decisão adequada à escala do acervo (centenas de estabelecimentos), em que a busca exata custa poucos milissegundos (Seção~\ref{sec:val_mlops}).

\subsection{Classificadores estatísticos clássicos para vetores esparsos}
Sobre as representações esparsas de alta dimensionalidade produzidas pelo TF-IDF, duas famílias de classificadores estatísticos clássicos mostram-se particularmente adequadas e computacionalmente econômicas, sendo adotadas como linhas de base na avaliação experimental (Capítulo~\ref{Cap:Ava_Exp}).

\subsubsection{Naive Bayes multinomial e complementar}
O classificador \textit{Naive Bayes} aplica o teorema de Bayes sob a hipótese simplificadora (\textit{naive}) de independência condicional entre os atributos dada a classe. Na variante Multinomial \cite{mccallum1998comparison}, cada atributo (termo do vocabulário) é modelado pela sua contagem/frequência de ocorrência no documento, estimando-se a probabilidade \textit{a posteriori} de uma classe $c$ dado um vetor de termos $x$ por:
\begin{equation}
	P(c \mid x) \propto P(c) \prod_{i=1}^{n} P(t_i \mid c)^{x_i}
\end{equation}
em que $P(c)$ é a probabilidade \textit{a priori} da classe, $t_i$ é o $i$-ésimo termo do vocabulário de tamanho $n$, $P(t_i \mid c)$ é a probabilidade de ocorrência desse termo na classe e o expoente $x_i$ é a contagem (ou peso TF-IDF) do termo $t_i$ no documento. Essa formulação é nativamente compatível com vetores TF-IDF esparsos e de alta dimensionalidade, pois a complexidade de treino e inferência escala linearmente com o número de termos não nulos, dispensando o cálculo de matrizes de covariância densas exigido por classificadores gaussianos. A variante Complementar (\textit{Complement Naive Bayes}) \cite{rennie2003tackling} corrige o viés sistemático do Multinomial em favor das classes majoritárias sob desbalanceamento, calculando as estatísticas de cada classe a partir do complemento (isto é, de todas as demais classes). Esta pesquisa avaliou a variante Multinomial com suavização de Laplace (Capítulo~\ref{Cap:Ava_Exp}); a variante Complementar, assim como as técnicas de reamostragem para desbalanceamento, não foi avaliada e é apontada como desdobramento imediato (Seção~\ref{sec:concl_limitacoes}).

\subsubsection{Florestas aleatórias (Random Forest)}
As Florestas Aleatórias \cite{breiman2001random} constituem um método de \textit{ensemble} que agrega as predições de múltiplas árvores de decisão treinadas independentemente sobre subamostras \textit{bootstrap} dos dados (\textit{bagging}) e sobre subconjuntos aleatórios de atributos sorteados em cada nó de divisão. A predição final é obtida por voto majoritário:
\begin{equation}
	\hat{y} = \text{moda}\{h_1(x), h_2(x), \dots, h_B(x)\}
\end{equation}
em que $h_b(x)$ denota a predição da $b$-ésima árvore, treinada sobre uma amostra \textit{bootstrap} distinta. A aleatoriedade na seleção de atributos por nó descorrelaciona as árvores individuais, reduzindo a variância do estimador agregado sem elevar substancialmente seu viés \cite{breiman2001random}. Florestas Aleatórias lidam nativamente com espaços de atributos esparsos e de alta dimensionalidade (como TF-IDF), são robustas a atributos irrelevantes e dispensam normalização prévia dos dados, o que as consolida como uma linha de base cientificamente relevante frente a arquiteturas neurais computacionalmente mais custosas.

\subsection{Métricas de avaliação para classificação multiclasse}
\label{sec:fund_metricas_classificacao}
O desempenho de um classificador multiclasse com $K$ categorias é formalmente descrito, seguindo a análise sistemática de medidas de \citeonline{sokolova2009systematic}, para cada categoria $k$, a partir de três contagens extraídas da matriz de confusão: Verdadeiros Positivos ($TP_k$), Falsos Positivos ($FP_k$) e Falsos Negativos ($FN_k$). A \textbf{Precisão} da categoria $k$ mede a proporção de predições daquela categoria que estavam corretas, e a \textbf{Revocação} (\textit{recall}) mede a proporção de instâncias verdadeiramente pertencentes à categoria $k$ que foram corretamente identificadas:
\begin{equation}
	\text{Precisão}_k = \frac{TP_k}{TP_k + FP_k}, \qquad
	\text{Revocação}_k = \frac{TP_k}{TP_k + FN_k}
\end{equation}
O \textbf{F1-Score} por categoria é a média harmônica entre Precisão e Revocação, penalizando desequilíbrios acentuados entre as duas grandezas:
\begin{equation}
	\text{F1}_k = \frac{2 \cdot \text{Precisão}_k \cdot \text{Revocação}_k}{\text{Precisão}_k + \text{Revocação}_k}
\end{equation}
Para consolidar essas métricas por categoria em um único indicador do classificador, duas estratégias de agregação são adotadas na literatura \cite{sokolova2009systematic}. A \textbf{média Macro} calcula a média aritmética simples das métricas de cada categoria, atribuindo peso igual a todas elas independentemente do seu volume de amostras:
\begin{equation}
	\text{F1-Macro} = \frac{1}{K} \sum_{k=1}^{K} \text{F1}_k
\end{equation}
Já a \textbf{média Ponderada} (\textit{weighted average}) pondera a métrica de cada categoria $k$ pelo seu suporte $n_k$ (número de instâncias verdadeiras daquela categoria) em relação ao total de amostras $N$:
\begin{equation}
	\label{eq:f1_ponderado}
	\text{F1-Ponderado} = \sum_{k=1}^{K} \frac{n_k}{N} \cdot \text{F1}_k
\end{equation}
A distinção entre as duas médias é decisiva sob desbalanceamento de classes: a média Ponderada tende a refletir principalmente o desempenho nas categorias majoritárias, ao passo que a média Macro penaliza igualmente o fraco desempenho em categorias minoritárias e de cauda longa, sendo por isso adotada como métrica primária nesta pesquisa (Capítulo~\ref{Cap:Ava_Exp}).

\section{Visão computacional e inteligência de documentos}
\label{sec:fund_visao_ocr}

A Inteligência de Documentos (\textit{Document AI}) integra Visão Computacional e PLN para extrair informações estruturadas de imagens e documentos não nativos \cite{docling2024}.

\subsection{Análise de leiaute de documentos (DLA) e caixas delimitadoras}
\label{sec:fund_dla}
A Análise de Leiaute de Documentos (DLA) é a etapa encarregada de segmentar a imagem da página, identificando blocos lógicos como tabelas, colunas, parágrafos e cabeçalhos -- problema clássico da área de análise de documentos, cujas abordagens evoluíram de heurísticas geométricas para detectores neurais de regiões \cite{binmakhashen2019document}. Cada entidade visual é mapeada por coordenadas espaciais denominadas caixas delimitadoras (\textit{bounding boxes}): $[x_{\min}, y_{\min}, x_{\max}, y_{\max}]$.

Em faturas de cartão de crédito, a preservação da topologia e da ordem de leitura é obrigatória: a relação espacial de alinhamento horizontal entre o nome do estabelecimento comercial e o seu respectivo valor monetário em reais (R\$) carrega a semântica relacional necessária para isolar a transação e evitar a fusão indevida de dados entre colunas paralelas \cite{docling2024}.

\subsection{Paradigmas de ingestão documental: reconhecimento óptico versus extração de fluxo nativo}
\label{sec:fund_paradigmas_ingestao}
A ingestão automatizada de documentos comerciais admite dois paradigmas conceitualmente distintos, cuja escolha depende da natureza do arquivo de origem.

O primeiro paradigma é o \textbf{Reconhecimento Óptico de Caracteres sobre imagem rasterizada} (\textit{raster-based OCR}), aplicável quando o documento de origem não possui nenhuma camada de texto digitalmente codificada -- como em faturas escaneadas ou fotografadas por câmeras móveis. Nesse cenário, o documento chega à esteira como uma matriz de pixels, e um modelo de reconhecimento (tipicamente redes neurais recorrentes do tipo LSTM, a exemplo do motor Tesseract \cite{smith2007overview}) infere, símbolo a símbolo, qual caractere melhor explica cada região da imagem. Trata-se de um problema inerentemente probabilístico: a qualidade da inferência degrada-se diante de ruídos visuais (baixa resolução, amassados, iluminação irregular), e mesmo em condições favoráveis persiste uma taxa de erro residual não nula.

O segundo paradigma é a \textbf{Análise de Leiaute com Extração de Fluxo Textual Nativo} (\textit{native text-stream extraction}), aplicável quando o documento de origem já contém uma camada de texto codificada digitalmente -- como em um PDF gerado programaticamente por um banco ou emissora de cartão. Nesse cenário, ferramentas de processamento documental (\textit{Document Parsing}), a exemplo do IBM Docling \cite{docling2024}, empregam modelos de visão computacional apenas para a Análise de Leiaute de Documentos (Seção~\ref{sec:fund_dla}) -- isto é, para segmentar tabelas, colunas e blocos textuais e inferir a ordem de leitura -- ao passo que os \textit{caracteres em si} são recuperados diretamente da codificação digital nativa do arquivo, sem qualquer etapa de reconhecimento óptico. Por não envolver inferência símbolo a símbolo sobre pixels, esse paradigma está, em princípio, livre do tipo de erro de transcrição mensurado pelo CER e pelo WER (Seção~\ref{sec:fund_metricas_ocr}) -- uma distinção retomada na discussão dos resultados do Capítulo~\ref{Cap:Ava_Exp}.

\subsection{Métricas formais de avaliação de OCR: CER e WER}
\label{sec:fund_metricas_ocr}
A qualidade da extração óptica é avaliada formalmente pela distância de edição de Levenshtein \cite{levenshtein1966binary} entre o texto predito pelo OCR e a verdade de referência. Definem-se três tipos fundamentais de erro: Substituição ($S$), Deleção ($D$) e Inserção ($I$).

A \textbf{Taxa de Erro de Caracteres} (CER -- \textit{Character Error Rate}) é calculada sobre o total de caracteres de referência $N_c$:
\begin{equation}
	\text{CER} = \frac{S_c + D_c + I_c}{N_c}
\end{equation}

A \textbf{Taxa de Erro de Palavras} (WER -- \textit{Word Error Rate}) opera de forma idêntica, porém adotando palavras inteiras como unidade amostral sobre o total de palavras $N_w$:
\begin{equation}
	\text{WER} = \frac{S_w + D_w + I_w}{N_w}
\end{equation}
Um único caractere lido incorretamente invalida a palavra inteira no cálculo do WER, tornando esta métrica estritamente mais rigorosa que o CER \cite{araujo2024saneamento}.

\section{Arquiteturas neurais e agentes inteligentes}
\label{sec:fund_arquiteturas_agentes}

\subsection{Mecanismos de autoatenção e arquiteturas Transformer}
A introdução do mecanismo de autoatenção (\textit{self-attention}) na arquitetura Transformer por Vaswani et al. (\citeyear{vaswani2017attention}) superou as limitações de paralelização das redes recorrentes tradicionais (LSTM e GRU) \cite{hochreiter1997long, cho2014learning}. Seu núcleo é a atenção escalada por produto escalar entre matrizes de consulta, chave e valor, que permite a cada \textit{token} ponderar dinamicamente todos os demais elementos da sequência de entrada e capturar dependências semânticas não locais mesmo em textos curtos e informais. Esta pesquisa não modifica, analisa nem mede o mecanismo de atenção: os modelos baseados em Transformer aqui empregados -- o codificador de sentenças MiniLM (Seção~\ref{sec:fund_pln}) e o ByT5 -- são usados como caixas-pretas pré-treinadas, razão pela qual sua formalização matemática (Vaswani et al., \citeyear{vaswani2017attention}) não é reproduzida. O que é instrumental a esta pesquisa é a forma como esses modelos tokenizam a entrada -- em subpalavras, no caso do MiniLM, e em bytes, no caso do ByT5 (Seção~\ref{sec:fund_tokenizacao}) --, pois é a tokenização, e não a atenção, que determina a robustez a nomes truncados e aglutinados.

\subsection{Grandes modelos de linguagem (LLMs) e raciocínio estruturado}
Os Grandes Modelos de Linguagem (\textit{Large Language Models} -- LLMs) pré-treinados em escala massiva demonstraram notável capacidade de inferência semântica e raciocínio contextual, inclusive em regime de poucos ou nenhum exemplo (\textit{few-shot} e \textit{zero-shot}), em que a tarefa é especificada apenas em linguagem natural no \textit{prompt} \cite{brown2020language}. Para mitigar o risco de alucinações e favorecer a previsibilidade das saídas exigida por sistemas financeiros, três técnicas são empregadas nesta pesquisa:
\begin{itemize}
	\item \textbf{Cadeia de Pensamento (\textit{Chain-of-Thought} -- CoT):} Técnica de indução por exemplos (\textit{prompting}) que leva o modelo a explicitar seus passos intermediários de raciocínio antes de emitir o veredito final, elevando a precisão lógica em tarefas que exigem inferência multietapa \cite{wei2022chain};
	\item \textbf{Chamada de Ferramentas (\textit{Tool Calling}) e Raciocínio-e-Ação:} Capacidade do modelo de suspender temporariamente a geração de texto para invocar funções externas (como motores de busca na web) a fim de obter evidências factuais em tempo real, intercalando passos de raciocínio verbalizado com ações concretas no padrão \textit{ReAct} \cite{yao2023react};
	\item \textbf{Saída Estruturada com Validação de Tipos:} Utilização de bibliotecas de validação de esquemas -- nesta pesquisa, a Pydantic, que valida dados a partir de anotações de tipo da linguagem Python \cite{pydantic2025docs} -- para forçar o modelo a responder exclusivamente em formatos JSON estritos e tipados.
\end{itemize}

\subsection{Orquestração baseada em grafos de estado: LangGraph}
Enquanto cadeias lineares simples (\textit{chains}) são rígidas, aplicações complexas demandam arquiteturas cíclicas e com controle fino de fluxo. O conceito de orquestração por grafo de estados remonta à formalização dos \textit{statecharts} \cite{harel1987statecharts}, que estende diagramas de transição de estados convencionais com hierarquia, concorrência e comunicação para especificar sistemas reativos complexos. O LangGraph \cite{langchain2024langgraph} aplica esse princípio à orquestração de agentes de IA, estruturando a aplicação como um Grafo Direcionado de Estados (StateGraph), em que cada nó executa uma função atômica (normalização, consulta vetorial, inferência LLM, persistência) e as arestas condicionais determinam o próximo passo baseando-se no estado tipado (\texttt{AgentState}). O arcabouço suporta nativamente persistência de memória (\textit{checkpoints}) e o padrão HITL, permitindo interromper o grafo de forma transparente para aguardar a confirmação humana e retomá-lo posteriormente.

\section{Engenharia de dados, MLOps e sustentabilidade de modelos}
\label{sec:fund_mlops}

A sustentabilidade de uma solução de IA em ambientes de produção exige a incorporação de práticas de Engenharia de Dados e Operações de Aprendizado de Máquina (\textit{Machine Learning Operations} -- MLOps) \cite{kreuzberger2022mlops, treveil2020mlops}.

\subsection{Desvios estatísticos: \textit{data drift} e \textit{concept drift}}
\label{sec:fund_drift}
Modelos de aprendizado de máquina operando sobre dados financeiros sofrem inevitável degradação temporal decorrente de dois fenômenos \cite{quinonero2008dataset, Lu2019LearningConceptDriftReview}. Registra-se desde já que esses fenômenos cumprem, nesta pesquisa, papel \textbf{motivador} do ciclo de MLOps e do aprendizado ativo (Seções~\ref{sec:sol_arquitetura} e~\ref{sec:val_mlops}), e não de objeto de medição: sua caracterização empírica sobre o acervo é apontada como trabalho futuro (Seção~\ref{sec:concl_limitacoes}).
\begin{itemize}
	\item \textbf{\textit{Data Drift} (Desvio de Covariáveis):} Ocorre quando a distribuição das variáveis de entrada $P(X)$ se altera ao longo do tempo (ex.: novos nomes de estabelecimentos no mercado, atualização de \textit{gateways} ou mudanças de formatação no OCR), embora a relação mapeada $P(Y|X)$ permaneça estável;
	\item \textbf{\textit{Concept Drift} (Desvio de Conceito):} Manifesta-se quando a probabilidade condicional real $P(Y|X)$ muda (ex.: uma grande rede de postos de combustível que passa a atuar predominantemente como loja de conveniência/alimentação, alterando a categoria esperada para uma mesma descrição textual).
\end{itemize}

\subsection{Aprendizado ativo (\textit{active learning}) por amostragem de incerteza}
Para contornar o custo proibitivo da rotulação manual exaustiva diante do surgimento contínuo de novas empresas, a literatura preconiza o Aprendizado Ativo (\textit{Active Learning}) \cite{ren2021aprendizagem}. Em vez de amostrar dados aleatoriamente, o sistema calcula a incerteza estatística de suas inferências -- a estratégia de amostragem por incerteza introduzida por \citeonline{lewis1994sequential} para o treinamento de classificadores de texto. Transações com escore de confiança inferior a um limiar seguro (nesta pesquisa, confiança $\le$~0,80, Seção~\ref{sec:sol_agente_inteligente}) são isoladas e encaminhadas para curadoria humana (HITL) -- a estratégia de amostragem por incerteza. Os dados confirmados ampliam seletivamente a base de conhecimento. Um problema correlato, mas distinto, é o do \textbf{ruído nos rótulos já existentes}: métodos como o \textit{Confident Learning} \cite{northcutt2021aprendizagem} estimam a distribuição conjunta entre rótulos observados e rótulos verdadeiros para identificar exemplos provavelmente mal rotulados. Esse método não é empregado nesta pesquisa; sua pertinência é discutida na Seção~\ref{sec:val_circularidade}, a propósito do risco de circularidade entre a rotulagem semiautomática e o desempenho medido.

\subsection{Arquitetura de \textit{data lakehouse} e rastreamento com MLflow}
A infraestrutura de suporte a dados moderna adota o paradigma de \textit{data lakehouse} \cite{armbrust2021lakehouse}, combinando a flexibilidade e o baixo custo de armazenamento de objetos compatíveis com S3 (como o MinIO) com as garantias transacionais e a organização em camadas Bronze/Silver/Gold -- a chamada \textit{Medallion Architecture}, convenção de mercado difundida pela Databricks a partir de sua documentação técnica, e não uma contribuição da literatura acadêmica \cite{armbrust2021lakehouse, treveil2020mlops}:
\begin{itemize}
	\item \textbf{Camada Bronze:} Armazena os dados brutos e imutáveis exatamente como foram ingeridos (PDFs originais, extratos bancários em CSV, textos ruidosos extraídos);
	\item \textbf{Camada Silver:} Contém os dados limpos, saneados, normalizados e enriquecidos, organizados em esquemas analíticos;
	\item \textbf{Camada Gold:} Tabelas agregadas e visões analíticas prontas para consumo por motores de consulta e ferramentas de inteligência de negócios.
\end{itemize}

O ciclo de vida experimental e operacional é unificado pelo MLflow \cite{zaharia2018mlflow, treveil2020mlops}. A plataforma atua como repositório centralizado de parâmetros, métricas e artefatos, suportando execuções aninhadas (\textit{nested runs}) que viabilizam o rastreamento em dois níveis de granularidade: a corrida global correspondente à fatura completa e corridas filhas isoladas para cada transação individualmente classificada.

\section{Considerações finais do capítulo}
\label{sec:fund_consideracoes_finais}

Os conceitos apresentados neste capítulo constituem a fundamentação técnica indispensável para a construção da solução proposta. A integração entre normalização morfológica, \textit{embeddings} densos no PostgreSQL com pgvector, análise de leiaute com Docling, orquestração em grafos com LangGraph e práticas contínuas de MLOps no \textit{data lakehouse} estabelece um arcabouço metodológico sólido para o desafio de classificar faturas de cartão de crédito sob restrições de baixo custo computacional -- cuja eficácia preditiva é apurada empiricamente nos Capítulos~\ref{Cap:Ava_Exp} e~\ref{Cap:Conclusoes}, e não pressuposta aqui. O capítulo seguinte revisa os trabalhos correlatos que empregam esses fundamentos, posicionando cientificamente as escolhas arquiteturais desta pesquisa frente ao estado da arte.
